% Appendix F -- Additional results (online appendix; draft, 30 October 2026)
%   tables    replication/output/tables/tab-mechanisms-payoff.tex
%             (code/4-exhibits/tab_mechanisms.R): Table 7 of Section 5.1 until
%             30 October, replaced there by fig:mechanisms-payoff
%             (code/4-exhibits/fig_mechanisms_payoff.R)
\section{Additional Results}\label{app:results-more}

Table~\ref{tab:mechanisms-payoff} gives the numbers behind
Figure~\ref{fig:mechanisms-payoff}, with the efficient share at 35--64,
which the figure leaves out, and purchases at 65--75. The share at 35--64
lies between the shares at 25--34 and 65--75 in every row but one: at
$\varrho=0.96$ it is the highest of the three, 45.9 percent. Purchases at
65--75 rise with comfort, from 5.5 per 100 owners a year without it to 6.6
with comfort doubled, but stay at 6.0 to 6.1 as $\varrho$ rises.

\begin{table}[htbp]
  \centering
  \caption[The efficient share of owners when the payoff changes]{The efficient share of owners when the payoff changes}\label{tab:mechanisms-payoff}
  \begin{threeparttable}
    \small
    \pkgtable{tab-mechanisms-payoff}
    \begin{tablenotes}[flushleft]\footnotesize
      \item \textit{Notes:} The runs of
      Figure~\ref{fig:mechanisms-payoff}. Each row changes one element of the
      model, keeps the other parameters at the estimates of
      Table~\ref{tab:estimates} (values in parentheses), and lets households
      re-optimize. Shares are A--C owners as a percentage of owners in each
      age group, from 65,536 simulated households. The gap is the share at
      25--34 less the share at 65--75, in percentage points; the ratio is
      the share at 25--34 over the share at 65--75. EHS: the targets of
      Table~\ref{tab:estimates}.
      \item $^a$All home purchases, by renters and owners, per 100
      owners aged 65--75 a year.
      \item $^b$The rows for $\varrho$ are simulated on a CPU and the
      others on a GPU. The CPU simulation at the estimates gives shares of
      22.5, 17.7, and 12.7 percent, a gap of 9.9, and a ratio of 1.78; the
      two differ by simulation noise only.
    \end{tablenotes}
  \end{threeparttable}
\end{table}
